\begin{table}[t]
\centering
\caption{The assignment family with a frozen and a fine-tuned backbone (test R@1, \%, mean$\pm$std over 3 seeds; hard: mean R@1 over C2, C5, C6). All runs in this table were made on the same GPUs. The OT rows share the aggregation architecture and differ only in the assignment; NetVLAD aggregates residuals to $k$-means-initialised centroids and has no dust-bin.}
\label{tab:family}
\footnotesize
\setlength{\tabcolsep}{3pt}
\begin{tabular}{@{}lcccc@{}}
\toprule
& \multicolumn{2}{c}{Frozen} & \multicolumn{2}{c}{Fine-tuned} \\
\cmidrule(lr){2-3}\cmidrule(lr){4-5}
Assignment & R@1 & hard & R@1 & hard \\
\midrule
$\lambda=0$, dust-bin & 60.1{\scriptsize$\pm$0.9} & 48.1 & 73.1{\scriptsize$\pm$1.1} & 61.9 \\
$\lambda=0$, no dust-bin & 60.3{\scriptsize$\pm$1.0} & 48.9 & 71.9{\scriptsize$\pm$1.4} & 61.2 \\
$\lambda=0.25$ & 60.0{\scriptsize$\pm$1.1} & 48.1 & 73.0{\scriptsize$\pm$1.4} & 62.4 \\
$\lambda=0.5$ & 59.4{\scriptsize$\pm$0.9} & 47.3 & 72.0{\scriptsize$\pm$0.8} & 60.7 \\
learned $\lambda$ (Galatea, init.\ 0.5) & 59.5{\scriptsize$\pm$1.3} & 47.4 & 71.3{\scriptsize$\pm$1.2} & 60.1 \\
$\lambda=0.75$ & 58.6{\scriptsize$\pm$1.0} & 47.2 & 73.3{\scriptsize$\pm$0.1} & 61.6 \\
$\lambda=1$ (SALAD), dust-bin & 59.6{\scriptsize$\pm$0.4} & 48.7 & 73.5{\scriptsize$\pm$0.5} & 62.2 \\
$\lambda=1$, no dust-bin & 55.8{\scriptsize$\pm$1.5} & 44.1 & 74.4{\scriptsize$\pm$0.9} & 62.9 \\
\midrule
NetVLAD (residual, $k$-means) & 57.8{\scriptsize$\pm$0.7} & 46.3 & 76.4{\scriptsize$\pm$0.7} & 64.3 \\
\bottomrule
\end{tabular}
\end{table}

\subsection{The assignment family with a frozen backbone}\label{sec:family}\label{sec:stage1}
Training only the aggregation head on illumination-diverse LunarSim scenes lifts R@1 from
$\zsSaladRone\%$ to $\galRone\%$ (Galatea head, three seeds; paired gain $\gainZsMean$ points, 95\% CI
$[\gainZsLo,\gainZsHi]$) and mean R@1 over the three hardest conditions (C2, C5, C6) from
$\zsSaladHard\%$ to $\galHard\%$ (\cref{tab:main}). All cluster-based heads clearly beat every
zero-shot model, whereas GeM pooling of the same frozen tokens does not ($\gemRone\%$): global pooling
cannot separate place-specific from illumination-dependent tokens~\cite{izquierdo2024optimal}.

\Cref{tab:family} (left) confirms predictions (P1) and (P2). Every member of the OT family, from
$\lambda=0$ to $\lambda=1$, reaches the same R@1 within seed noise (spread across fixed $\lambda$:
$\famFrLamRange$ points), and Galatea's learned head is statistically indistinguishable from SALAD
($\galVsSalMean$ points, CI $[\galVsSalLo,\galVsSalHi]$). A learned $\lambda$ barely leaves its
initialisation ($\bar\lambda=\galLam$); even with a $25\times$ larger learning rate on $\lambda$ and
$z$ it only reaches $\galFastLam$, at unchanged accuracy ($\galFastRone\%$). This happens although the
heads spend mass very differently: the balanced head sends $\salDust\%$ of the token mass to the
dust-bin, an amount that tracks the sky fraction of the image ($\rho=\salSkyCorr$), the learned head
only $\galDust\%$ (coefficient of variation of cluster mass $\galCV$ vs.\ $\salCV$). As
\cref{prop:colinv} states, intra-normalisation removes exactly this difference.

The dust-bin shows where the theory's assumptions bite. At $\lambda=0$ it changes only the token
weights $w_i$ and is irrelevant ($\famFrNDvsDMean$ points without it, CI $[\famFrNDvsDLo,\famFrNDvsDHi]$).
At $\lambda=1$ it is essential ($\famFrSalNDvsDMean$, CI $[\famFrSalNDvsDLo,\famFrSalNDvsDHi]$): without
a place to put surplus mass, the balanced constraint forces every token, sky included, into the
clusters in fixed proportions, which no re-weighting can undo. Residual NetVLAD aggregation is slightly
below the OT heads on frozen tokens ($\famFrNv\%$).

\subsection{The assignment family with a fine-tuned backbone}\label{sec:ft}
Fine-tuning the last four backbone blocks together with the head (\cref{tab:main}, bottom) improves
every head substantially: GeM by $\ftGemGain$, NetVLAD by $\ftNvGain$, SALAD by $\ftSalGain$ and the
Galatea head by $\ftGalGain$ points R@1 (CI $[\ftGainLo,\ftGainHi]$; hard conditions:
$\galHard\%\to\ftGalHard\%$). The family no longer collapses completely (\cref{tab:family}, right),
but its structure is simple. All \emph{fixed} values of $\lambda$ remain equivalent
(spread $\famFtLamRange$ points), and the dust-bin changes R@1 by about one point at most
($\lambda=0$: $\famFtNDvsDMean$ $[\famFtNDvsDLo,\famFtNDvsDHi]$; $\lambda=1$:
$\famFtSalNDvsDMean$ $[\famFtSalNDvsDLo,\famFtSalNDvsDHi]$), presumably because the backbone learns to
make uninformative tokens harmless. The \emph{learned} $\lambda$, which again stays near its
initialisation ($\bar\lambda=\ftGalLam$), is the one member that falls behind
($\famFtGal\%$; $\famFtGalVsSalMean$, CI $[\famFtGalVsSalLo,\famFtGalVsSalHi]$ vs.\ SALAD). LAMHALF

The largest difference is not inside the OT family at all. NetVLAD, which aggregates residuals to
$k$-means-initialised centroids, becomes the best head ($\ftNvRone\pm\ftNvRonesd\%$). Its closest
OT relative, $\lambda=0$ without a dust-bin, uses the same softmax assignment and no dust-bin, yet is
$\famFtNvVsNDMean$ points behind (CI $[\famFtNvVsNDLo,\famFtNvVsNDHi]$). With trainable tokens the
\emph{aggregation}, not the assignment, is therefore the design axis that matters; on frozen tokens the
same comparison favours the OT head ($\famFrNvVsNDMean$ points for NetVLAD).
